% !TeX program = lualatex
% Compile twice: lualatex 157.tex
% All references and prose are embedded; no BibTeX database or external inputs.
\documentclass[11pt,a4paper]{article}
\usepackage[margin=24mm,headheight=14pt]{geometry}
\usepackage{fontspec}
\setmainfont{Latin Modern Roman}
\setsansfont{Latin Modern Sans}
\setmonofont{Latin Modern Mono}
\usepackage{amsmath,amssymb,mathtools}
\usepackage{unicode-math}
\setmathfont{Latin Modern Math}
\usepackage{microtype}
\usepackage{enumitem,xurl,needspace}
\usepackage[dvipsnames]{xcolor}
\usepackage{hyperref}
\hypersetup{unicode=true,colorlinks=true,linkcolor=MidnightBlue,urlcolor=MidnightBlue,
 pdftitle={500 Mathematical Problem Statements},pdfauthor={Source-annotated compilation},bookmarksnumbered=true}
\usepackage{bookmark}
\usepackage{tocloft}
\setlength{\cftsecnumwidth}{3em}
\cftsetpnumwidth{2.5em}
\cftsetrmarg{4em}
\usepackage{titlesec}
\titleformat{\section}{\Large\bfseries\raggedright}{\thesection}{0.7em}{}
\usepackage{fancyhdr}
\pagestyle{fancy}\fancyhf{}
\fancyhead[L]{\small\sffamily Mathematical problem statements}
\fancyhead[R]{\small Source edition 22}
\fancyfoot[C]{\thepage}
\setlength{\parindent}{0pt}
\setlength{\parskip}{5pt plus 1pt}
\setlength{\emergencystretch}{3em}
\setcounter{tocdepth}{1}
\setcounter{secnumdepth}{1}
\setlist[itemize]{leftmargin=3em,itemsep=5pt,topsep=4pt,parsep=0pt}
\allowdisplaybreaks
\newcommand{\N}{\mathbb{N}}
\newcommand{\Z}{\mathbb{Z}}
\newcommand{\Q}{\mathbb{Q}}
\newcommand{\R}{\mathbb{R}}
\newcommand{\C}{\mathbb{C}}
\newcommand{\F}{\mathbb{F}}
\newcommand{\A}{\mathbb{A}}
\newcommand{\PP}{\mathbb P}
\newcommand{\E}{\mathbb{E}}
\newcommand{\eps}{\varepsilon}
\newcommand{\dd}{\,\mathrm d}
\newcommand{\sref}[2]{\hyperlink{source:#1:#2}{[S#2]}}
\newcommand{\eref}[2]{\hyperlink{extra:#1:#2}{[E#2]}}
% Turn the next switch off for a shorter reading copy. The quotations preserve
% every exactTarget field, including qualifications omitted from a short summary.
\newif\ifshowcatalogue
\showcataloguetrue
\newcommand{\cataloguescope}[1]{\ifshowcatalogue
 \paragraph{Exact scope in the supplied catalogue.}
 \begin{quote}\small #1\end{quote}\fi}

% Standalone research-notebook wrapper; the source preamble above is retained.
\numberwithin{equation}{section}
\hypersetup{pdftitle={Research attempt -- problem 157},
 pdfauthor={Research notebook; original compilation retained}}
\fancyhead[L]{\small\sffamily Research notebook}
\fancyhead[R]{\small\sffamily Problem 157 | 27 September 2026}
\begin{document}
\setcounter{section}{156}
% ================================================================
% problemId: problem.grothendieck-variational-hodge-conjecture
% source releaseRank: 157; source releaseStatus: open_with_solved_subcases
% exactTarget SHA-256: f403dcf4474779211445e8c85b2cbaa3399b6dc2015860b416a7a0f1feb3954b
\clearpage
\section{\texorpdfstring{Grothendieck's Variational Hodge Conjecture}{Grothendieck's Variational Hodge Conjecture}}\label{157}
\subsection{Definitions and mathematical statement}
For a smooth projective family $f:X\to S$, the relative de Rham bundle
$\mathcal H^{2p}_{\mathrm{dR}}=\mathbf R^{2p}f_*\Omega^\bullet_{X/S}$ has the Gauss--Manin connection $\nabla$. A global algebraic section $\xi$ is horizontal when $\nabla\xi=0$, so analytically it is transported by the cohomological local system. Let $\mathrm{cl}_{\mathrm{dR}}$ be the cycle-class map from rational codimension-$p$ cycles on a fiber. The target is
\[
 \xi(s_0)\in\operatorname{im}\mathrm{cl}_{\mathrm{dR},s_0}
 \quad\Longrightarrow\quad
 \xi(s)\in\operatorname{im}\mathrm{cl}_{\mathrm{dR},s}\quad\text{for all }s\in S(\C).
\]
Here $S$ is smooth and connected. No single relative cycle realizing all fibers is demanded; only fiberwise algebraicity of the transported class is asserted.
\par\smallskip\textit{Formulation and concept references:} \sref{157}{1}.
\cataloguescope{For every smooth projective morphism f:X->S of smooth connected complex algebraic varieties, every integer p>=0, and every global algebraic horizontal section xi of the relative algebraic de Rham cohomology bundle R\textasciicircum{}(2p)f\_*Omega\textasciicircum{}.\_(X/S) for the Gauss-Manin connection, if xi(s0) is the de Rham cycle class of a rational algebraic codimension-p cycle on one fiber X\_s0, then xi(s) is such a rational algebraic cycle class on every fiber X\_s for s in S(C).}
\subsection{Short English statement}
If an algebraic cohomology class is transported horizontally through a smooth projective family, must it remain algebraic on every fiber?
% BEGIN RESEARCH ADDITION ATTEMPT-157
\subsection{Research attempt}

\paragraph{Current frontier.}
The exact global formulation matters. Charles--Schnell's Theorem 33 and
Proposition 34 show that a global flat section which is a rational Hodge
class at one point remains Hodge everywhere, and that the Hodge conjecture
would imply this variational statement \sref{157}{1}. Thus an analytic
disk on which a locally transported class leaves its Hodge locus is not a
counterexample to the printed global algebraic formulation. Recent relevant
work includes Pridham's 2024 revision on semiregularity and obstruction
maps \eref{157}{1}, and Bae--Kool--Park's 2026 publication giving a
variational Hodge application for Calabi--Yau fourfold families with a
nonzero reduced virtual cycle \eref{157}{3}. Such extra geometric or
nonvanishing conditions are not available for an arbitrary input here.

\paragraph{Attack route.}
The Hodge-conjecture route imports a larger unresolved statement; deformation
theory must remove actual cycle obstructions, not merely preserve Hodge type.
We instead parameterize representatives by proper Hilbert schemes and ask
whether one parameter component dominates the base. We then test whether
first-order semiregularity can supply this missing dominance.

\paragraph{Attempt: rationality and the correct local reduction.}
Via the comparison isomorphism, with the Tate twist consistent with the
cycle-class normalization, $\xi$ corresponds to a global section
\[
 \alpha\in H^0(S^{\mathrm{an}},R^{2p}f_*\Q(p)).
\]
Indeed, flatness gives a complex local-system section and rationality at
$s_0$ propagates along paths. The fixed-part conclusion from the cited
invariant-cycle theorem gives the Hodge condition throughout the base
\sref{157}{1}, Propositions 32 and 34.

We first work over a nonempty affine open in $S$, so that projective
embeddings and the quoted quasi-projective theorems apply. This is not a
permanent restriction of the target. A smooth connected complex algebraic
variety is irreducible, so any two nonempty affine opens intersect. A
universal result on affine bases, started on one containing $s_0$, would
therefore propagate across an affine cover by taking an algebraic fiber
class at an overlap as the next starting point. The relative dimension
$n$ is constant. For $p>n$ the target class is zero; below assume
$0\le p\le n$.

\paragraph{Attempt: a countable proper parameterization.}
Fix a relatively very ample line bundle. For finite data
\[
 m\ge1,\quad r\ge1,\quad a_1,\ldots,a_r\in\Z,
 \quad P_1,\ldots,P_r\text{ Hilbert polynomials of degree }n-p,
\]
consider
\begin{equation}\label{eq157-hilbert}
 T=\prod_{i=1}^r{}_S\operatorname{Hilb}_{P_i}(X/S).
\end{equation}
This is a projective finite-type $S$-scheme, with universal flat closed
subschemes $Z_i\subset X\times_ST$; see the Hilbert case of
\eref{157}{2}, Theorem 5.1. Take their fundamental cycles of relative
dimension $n-p$, ignoring lower-dimensional embedded contributions.
They determine a section
\[
 \sum_i a_i\operatorname{cl}([Z_{i,t}])-m\alpha_{f_T(t)}
\]
of the pulled-back rational cohomology local system on $T^{\mathrm{an}}$.
Here we use the usual specialization compatibility of fundamental cycle
classes in a proper flat family: along a curve in the parameter space,
the specialized fundamental cycle has the transported cohomology class.
For an irreducible parameter space, curves through a point and a dense
open, together with local constancy on the smooth locus, give the same
conclusion across its singular points. This explains why flatness of the
\emph{universal subscheme}, rather than smoothness of the Hilbert scheme,
is the relevant input.

Restrict to each reduced irreducible component $T_j$ of $T$ on which that
section vanishes at some point. It then vanishes identically, since an
irreducible complex algebraic variety is analytically path connected.
Let $A_j\subseteq S$ be the image of $T_j$. Properness makes $A_j$ Zariski
closed. There are only countably many choices of the finite integer and
polynomial data, and finitely many components for each choice. Thus the
algebraicity locus of this one class has the description
\begin{equation}\label{eq157-locus}
 \mathcal A_\alpha:=
 \{s:\alpha_s\text{ is a rational algebraic cycle class}\}
 =\bigcup_{j\ge1}A_j.
\end{equation}

Both inclusions deserve checking. A point of $A_j(\C)$ has a complex
preimage, whose weighted universal cycles divided by $m$ represent
$\alpha_s$. Conversely, an arbitrary rational cycle is a \emph{finite}
integer linear combination of integral subvarieties after multiplying by
some $m$. Those subvarieties define points of the corresponding Hilbert
schemes, and the containing component gives one of the $A_j$ above.
Negative coefficients and varying denominators are indispensable; using
only effective cycles of one fixed degree would not parameterize the
stated problem.

\textbf{Baire--properness lemma.} If $\mathcal A_\alpha$ is not meagre in
$S^{\mathrm{an}}$, then $A_j=S$ for some $j$, and in particular
$\mathcal A_\alpha=S(\C)$.

\textit{Proof.} A proper algebraic closed subset of the smooth irreducible
base has empty analytic interior and is nowhere dense. If every $A_j$ were
proper, \eqref{eq157-locus} would be meagre. This contradicts the
hypothesis. Nonempty analytic open sets are not meagre, since complex
manifolds are locally compact Baire spaces. \hfill$\square$

Increasing finite unions of the $A_j$ are closed, but their union need not
be. Membership of $s_0$ supplies no compactness argument for dominance.

\paragraph{Attempt: trace from a dominating parameter space.}
There is a stronger useful consequence of $A_j=S$ that can be proved rather
than added to the target. Choose a closed point of the generic fiber of
$T_j\to S$ and take its reduced closure $T'\subseteq T_j$. The residue
field of this point is finite over $\C(S)$, so $T'\to S$ is proper,
dominant, and generically finite, of some degree $e\ge1$. After shrinking
to a nonempty open $U\subset S$, it is finite \emph{\'etale}, since the
field extension is separable in characteristic zero. Pull back the
universal subschemes to $T'$. Their flatness is preserved by base change.

Let $q:X\times_ST'\to X$ be the proper projection. Taking fundamental
cycles and proper pushforwards defines
\begin{equation}\label{eq157-trace}
 Z=\frac1{em}\,q_*\left(\sum_i a_i[Z_i|_{T'}]\right)
       \in\mathrm{CH}^p(X)_{\Q}.
\end{equation}
The codimension is correct because $T'$ has the same dimension as $S$.
For $s\in U(\C)$, the fiber has $e$ preimages, and specialization and
proper pushforward give
\[
 \operatorname{cl}(Z|_{X_s})
 =\frac1{em}\sum_{t\in T'_s}\sum_i a_i\operatorname{cl}(Z_{i,t})
 =\alpha_s.
\]
The restriction of a cohomology class on $X$ to the smooth fibers is a
global flat section. It follows by connectedness that equality holds on
\emph{all} of $S$. The restriction of the Chow class is interpreted by
refined Gysin pullback along $X_s\hookrightarrow X$, a regular embedding
because $f$ and $S$ are smooth; flatness of the pushed-forward cycle over
all of $S$ is not required.

Thus dominance gives a global rational cycle with the correct fiber
\emph{classes}. This deduction does not assert that a specified cycle at
$s_0$ deforms, that one effective flat family works everywhere, or that an
integral cycle is obtained. Division by $em$ is legitimate precisely
because the target uses rational coefficients.

\paragraph{Attempt: a precise sufficient deformation input.}
Consider the following proposed input for each family and class:
\begin{quote}
After replacing the initial rational cycle by a rationally weighted tuple
with the same cohomology class, there are data as in
\eqref{eq157-hilbert} and a point $t_0$ over $s_0$, with
$\sum_i a_i[Z_{i,t_0}]=m\alpha_{s_0}$ in cohomology, such that
$T\to S$ is smooth at $t_0$.
\end{quote}
If this input holds, the smooth locus containing $t_0$ maps onto a
nonempty open subset of $S$. Every sufficiently near parameter point has
the same class identity on its relevant component. Its images are
contained in $\mathcal A_\alpha$, so the Baire--properness lemma gives
$A_j=S$, and \eqref{eq157-trace} completes the fiberwise conclusion. This
is a conditional implication, not a proof of the quoted input.
Formal smoothness for the finite-presentation morphism would suffice;
first-order tangent surjectivity would not.

For an lci subvariety $Z\subset X_s$ of codimension $p$, the customary
semiregularity map has the form
\[
 H^1(Z,N_{Z/X_s})\longrightarrow
 H^{p+1}(X_s,\Omega_{X_s}^{p-1}).
\]
Pridham explains the annihilation of obstruction classes and its
ambient-deformation interpretation, with the more natural truncated
de Rham targets needed in full generality \eref{157}{1}, Introduction
and Sections 2.4--2.5. That information does not say that the kernel is
zero. For a signed tuple, even cancellation of the \emph{sum} of the
cohomological obstruction images need not kill each obstruction to
lifting the tuple. The missing bridge is therefore an all-orders
unobstructed representative, or another reason that one parameter
component dominates. It is not the already available persistence of
Hodge type.

\paragraph{Stress test.}
A concrete algebraic warning is the finite morphism
\[
 \operatorname{Spec}\C[t]/(t^2)\longrightarrow\A^1_{\C},
\]
induced by the coordinate $t$. At the origin its differential is
surjective between one-dimensional tangent spaces. The first-order base
arc $t=\varepsilon$ lifts over $\C[\varepsilon]/(\varepsilon^2)$, but the
same arc cannot lift over $\C[\varepsilon]/(\varepsilon^3)$ because its
square is nonzero. The image is just the origin. Thus even a proper map
with a surjective tangent map can fail completely to dominate the base.
This example refutes the proposed tangent-space shortcut, not VHC.

Nor does density of algebraic fibers suffice for the Baire lemma: a
countable union of proper closed subsets can be analytically dense.
Likewise, the presence of $s_0$ in \eqref{eq157-locus} gives only one
nonempty closed image. No reason has been established for it to contain
an open set. For divisors the Hodge condition plus the Lefschetz
$(1,1)$ theorem supplies fiberwise algebraicity, but using that special
case in higher codimension would import the Hodge conjecture itself
\sref{157}{1}, Proposition 34 and its discussion.

\paragraph{Outcome.}
\textbf{Outcome: Conditional reduction.}

The attempt proves a countable-proper-image description of the algebraicity
locus, a nonmeagreness-to-dominance criterion, and a trace construction that
turns a dominating cycle-parameter component into a global rational cycle
with the required fiber classes. This uses standard cycle and Hilbert-scheme
tools; no literature priority is claimed.

What remains is a dominating representative component starting from one
algebraic fiber, for example via all-orders smooth lifting. Semiregularity
alone has not supplied it. The affine-open overlap argument preserves the
full target but does not close this gap. No full proof or counterexample
is claimed.
% END RESEARCH ADDITION ATTEMPT-157
\subsection{Sources}
\begin{itemize}\raggedright
\item[\textnormal{[S1]}]\hypertarget{source:157:1}{} François Charles and Christian Schnell, Notes on absolute Hodge classes, Conjectures30/31, Proposition32, Theorem33 and Proposition34; base extension by the coordinator's recorded affine-open argument.
\url{https://arxiv.org/pdf/1101.3647}.
{\small\textit{Catalogue formulation source.}}
% BEGIN RESEARCH ADDITION SOURCES-157
\item[\textnormal{[E1]}]\hypertarget{extra:157:1}{} J. P. Pridham,
\emph{Semiregularity as a consequence of Goodwillie's theorem},
arXiv:1208.3111v4, 4 November 2024, Introduction and Sections 2.4--2.5.
\url{https://arxiv.org/html/1208.3111v4}.
{\small\textit{Consulted 27 September 2026; obstruction annihilation is not
used as an injectivity theorem for the semiregularity map.}}
\item[\textnormal{[E2]}]\hypertarget{extra:157:2}{} Nitin Nitsure,
\emph{Construction of Hilbert and Quot Schemes}, arXiv:math/0504590,
29 April 2005, Section 1 and Theorem 5.1 (Grothendieck).
\url{https://arxiv.org/html/math/0504590}.
{\small\textit{Consulted 27 September 2026; representability, universal
flat families, and projectivity for fixed Hilbert polynomial.}}
\item[\textnormal{[E3]}]\hypertarget{extra:157:3}{} Younghan Bae, Martijn Kool,
and Hyeonjun Park, \emph{Counting surfaces on Calabi--Yau 4-folds I:
Foundations}, \emph{Geometry \& Topology} 30 (2026), 2431--2564;
arXiv:2208.09474v2, 19 May 2025, abstract's variational Hodge application.
\url{https://arxiv.org/abs/2208.09474v2}.
{\small\textit{Consulted 27 September 2026. Used only for the stated
restricted frontier result, not as a theorem for arbitrary families.}}
% END RESEARCH ADDITION SOURCES-157
\end{itemize}

\end{document}
